% Generated draft: requires mathematical and source review.
\documentclass[11pt]{article}
\usepackage[margin=25mm]{geometry}
\usepackage{fontspec}
\setmainfont{Latin Modern Roman}
\usepackage{amsmath,amssymb,mathtools,unicode-math}
\setmathfont{Latin Modern Math}
\usepackage{enumitem,needspace,xurl,hyperref}
\setlength{\parindent}{0pt}
\setlength{\parskip}{5pt}
\setlength{\emergencystretch}{3em}
\newcommand{\N}{\mathbb{N}}
\newcommand{\Z}{\mathbb{Z}}
\newcommand{\Q}{\mathbb{Q}}
\newcommand{\R}{\mathbb{R}}
\newcommand{\C}{\mathbb{C}}
\newcommand{\F}{\mathbb{F}}
\newcommand{\A}{\mathbb{A}}
\newcommand{\PP}{\mathbb{P}}
\newcommand{\E}{\mathbb{E}}
\newcommand{\eps}{\varepsilon}
\newcommand{\dd}{\,\mathrm{d}}
\begin{document}
\section*{UnsolvedMath ID: 30002640}
\textbf{Source problem code:} catalogue entry
\subsection{Definitions and mathematical statement}
A Seifert space is a three-dimensional space equipped with a Seifert fibration by circles, with the usual local models allowing exceptional fibers. The supplied record does not specify the permitted boundary, compactness, or orientability assumptions, or whether the fibration is part of the data.

Write \(K=\Q(e^u)\) for the base field appearing in the source. The meaning of \(u\), and whether \(e^u\) is a formal indeterminate or an analytic quantity, are not specified. Let \(\mathcal I_A(M)\) denote the collection intended by the source phrase ``all genus-\(g\) \(A\)-series perturbative invariants of fiber knots in \(M\).'' This notation is a placeholder for a collection whose definition and index ranges require clarification.

The supplied finite-extension question has the following conditional mathematical form: for every Seifert space \(M\) satisfying the source condition \(\chi(M)>0\), does there exist a field \(L(M)\) containing \(K\) such that
\[
 [L(M):K]<\infty,
 \qquad
 \mathcal I_A(M)\subseteq L(M)?
\]
Here containment presupposes a specified interpretation of the invariants as field elements. The assertion asks for a single finite extension containing the entire intended collection for a given \(M\); separate algebraicity assertions for individual invariants would not express this uniform requirement.

The source does not define \(\chi(M)\). In particular, it must be clarified whether this denotes an Euler characteristic associated with the base orbifold of the Seifert fibration or the ordinary Euler characteristic of the underlying space. For a closed three-manifold the latter is zero, so that interpretation would make the stated positivity condition empty. No replacement of \(\chi(M)\) is made here.

The genus quantifier is also unresolved: the source does not say explicitly whether \(g\) is fixed or whether the same field must contain the invariants simultaneously for every \(g\in\Z_{\geq 0}\). Likewise, it does not define the \(A\)-series perturbative invariants, their remaining indices and normalizations, or the meaning of ``fiber knots.'' Consequently, the displayed statement records the field-theoretic requirement, but a fully specified conjecture requires these source conventions.
\subsection{Short English statement}
The question asks whether, for each eligible Seifert space, one finite extension of \(\Q(e^u)\) contains all the specified perturbative fiber-knot invariants. The meaning of the positivity condition, the invariants, and the genus quantifier remains unresolved in the supplied record. The supplied status is open; this draft does not independently verify that assessment.
\subsection{Sources}
\begin{itemize}
\item Local UnsolvedMath registry: supplied complete record for permanent ID \(30002640\), titled ``Algebraicity of Perturbative Fiber-Knot Invariants,'' source problem code \texttt{catalogue entry}, in the Oberwolfach Reports Open Problems collection. No individual proposer is supplied. The record gives proposed year 2014 and labels the statement \texttt{corrected\_verified}; these are supplied metadata, not independent mathematical verification.
\item Supplied report citation: \emph{Low-dimensional Topology and Number Theory} (2015), Oberwolfach Reports, pp.\ 2115--2176, DOI \texttt{10.4171/owr/2014/38}. \url{https://doi.org/10.4171/owr/2014/38}. The background describes these reports as documenting workshops and their open-problem discussions. The citation's year 2015 and the record's proposed year 2014 are retained as supplied.
\item Supplied literature reference: Borot--Eynard--Weisse, \url{https://doi.org/10.1007/s00029-016-0266-6}. No title, theorem number, or relevant passage is supplied.
\item Supplied literature reference labeled ``original OWR report'': \url{https://doi.org/10.4171/owr/2014/38}.
\item Literature assessment dated exactly \texttt{2026-08-21}: the supplied curation labels the problem open and reports that the finite-extension statement is proved for important spherical Seifert examples, while the general claim for all Seifert spaces and all stated perturbative invariants remains conjectural in the later primary treatment. This is supplied dated curation, not an independently verified assessment. The supplied background explicitly requests independent verification of time-sensitive status and cited claims.
\end{itemize}
\end{document}
